\documentclass[a4paper,11pt]{article}

% ── Pacotes ───────────────────────────────────────────────────────────────────
\usepackage[utf8]{inputenc}
\usepackage[T1]{fontenc}
\usepackage[brazil]{babel}
\usepackage[top=2.5cm, bottom=2.5cm, left=5.5cm, right=2cm]{geometry}
\usepackage{xcolor}
\usepackage{titlesec}
\usepackage{enumitem}
\usepackage{parskip}
\usepackage{fancyhdr}
\usepackage{eso-pic}
\usepackage{tikz}
\usepackage{mdframed}
\usepackage{hyperref}

% ── Cores ─────────────────────────────────────────────────────────────────────
\definecolor{azulescuro}{RGB}{30,55,100}
\definecolor{cinzaclaro}{RGB}{245,245,245}
\definecolor{laranja}{RGB}{210,100,30}

% ── Barra lateral esquerda ────────────────────────────────────────────────────
\AddToShipoutPictureBG{%
  \begin{tikzpicture}[remember picture, overlay]
    \fill[azulescuro]
      (current page.north west)
      rectangle
      ([xshift=4.8cm] current page.south west);
  \end{tikzpicture}%
}

% ── Cabeçalho / rodapé ───────────────────────────────────────────────────────
\pagestyle{fancy}
\fancyhf{}
\renewcommand{\headrulewidth}{0pt}
\fancyfoot[R]{\small\thepage}
\fancyfoot[L]{\small\color{gray}\today}

% ── Títulos de seção ─────────────────────────────────────────────────────────
\titleformat{\section}
  {\color{azulescuro}\large\bfseries}
  {}
  {0em}
  {}
  [\vspace{2pt}\color{laranja}\rule{\linewidth}{1.5pt}\vspace{4pt}]

% ── Estilo de listas ─────────────────────────────────────────────────────────
\setlist[itemize]{leftmargin=1.5em, itemsep=2pt, parsep=0pt}
\setlist[enumerate]{leftmargin=1.5em, itemsep=2pt, parsep=0pt}

% ── Comando: campo da barra lateral ──────────────────────────────────────────
% #1 = rótulo (maiúsculas)   #2 = valor
\newcommand{\campolateral}[2]{%
  {\color{white}\footnotesize\textbf{#1}}\\[2pt]%
  {\color{white!80!azulescuro}\footnotesize #2}\par\vspace{10pt}%
}

% =============================================================================
\begin{document}

% ── Barra lateral ─────────────────────────────────────────────────────────────
\begin{tikzpicture}[remember picture, overlay]
  \node[anchor=north west, xshift=0.35cm, yshift=-2.5cm,
        text width=4.1cm, align=left]
    at (current page.north west) {%
      \campolateral{NOME DO CANDIDATO:}{{{informacoes.nome_candidato}}}
      \campolateral{N\'UMERO DE INSCRI\c{C}\~AO:}{{{informacoes.numero_inscricao}}}
      \campolateral{\'AREA DO CONCURSO:}{{{informacoes.area_concurso}}}
      \campolateral{P\'UBLICO-ALVO:}{{{informacoes.publico_alvo}}}
      \campolateral{DISCIPLINA:}{{{informacoes.disciplina}}}
      \campolateral{TEMPO DA AULA:}{{{informacoes.tempo_aula}}}
      \campolateral{DATA DA AULA:}{{{informacoes.data_aula}}}
  };
\end{tikzpicture}

% ── Título principal ──────────────────────────────────────────────────────────
\vspace*{-1em}
{\fontsize{28}{32}\selectfont\bfseries PLANO DE AULA}

\vspace{0.3em}
{\color{laranja}\rule{0.08\linewidth}{2pt}}

\vspace{0.7em}
{\large\bfseries\color{azulescuro} {{titulo}} }

\vspace{1.2em}

% =============================================================================
\section{Objetivos}

% Backend: para cada entrada em objetivos[], gerar uma linha \item <texto>
\begin{itemize}
{{BLOCO_OBJETIVOS}}
\end{itemize}

% =============================================================================
\section{Conte\'udo}

% Backend: para cada objeto em conteudo[], gerar:
%   \item <titulo>
%   Se o objeto tiver subtopicos[], abrir enumerate aninhado com label=<N>.\arabic*.
%   e gerar um \item por subtópico, depois fechar o enumerate.
\begin{enumerate}
{{BLOCO_CONTEUDO}}
\end{enumerate}

% =============================================================================
\section{Metodologia de Ensino}

% Backend: para cada objeto em metodologia[], gerar:
%   \textbf{<nome>:} <descricao>\n\n
{{BLOCO_METODOLOGIA}}

% =============================================================================
\section{Recursos Did\'aticos}

% Backend: para cada entrada em recursos_didaticos[], gerar \item <texto>
\begin{itemize}
{{BLOCO_RECURSOS}}
\end{itemize}

% =============================================================================
\section{Atividade Avaliativa}

\textbf{Atividade proposta:} Ap\'os a exposi\c{c}\~ao, os estudantes receber\~ao uma
situa\c{c}\~ao-problema:

\begin{mdframed}[backgroundcolor=cinzaclaro, linecolor=azulescuro,
                 linewidth=1pt, innerleftmargin=10pt, innerrightmargin=10pt,
                 innertopmargin=8pt, innerbottommargin=8pt]
\textit{``{{atividade_avaliativa.enunciado}}''}
\end{mdframed}

\vspace{0.5em}
\textbf{Forma de aplica\c{c}\~ao:} {{atividade_avaliativa.forma_aplicacao}}

\vspace{0.4em}
\textbf{Crit\'erios de corre\c{c}\~ao:}

% Backend: para cada objeto em criterios_correcao[], gerar:
%   \item <criterio> — \textbf{<peso>\%}
\begin{itemize}
{{BLOCO_CRITERIOS}}
\end{itemize}

% =============================================================================
\section{Bibliografia de Refer\^encia}

% Backend: para cada objeto em bibliografia[], gerar uma linha \item formatada.
%   Se tiver campo "periodico"  -> formato de artigo:
%     AUTOR. \textbf{titulo}. \textit{periodico}, volume, ano.
%   Se tiver campo "editora"    -> formato de livro:
%     AUTOR. \textbf{titulo}. editora, ano.
\begin{itemize}[leftmargin=0pt, label={}]
{{BLOCO_BIBLIOGRAFIA}}
\end{itemize}

\end{document}